\documentclass[aspectratio=169, 10pt]{beamer}

% Paquetes esenciales para Ingeniería
\usepackage[utf8]{inputenc}
\usepackage[T1]{fontenc}
\usepackage[spanish]{babel}
\usepackage{amsmath, amssymb}
\usepackage{siunitx}
\usepackage{circuitikz}

% Configuración del tema Beamer (Estilo sobrio e institucional)
\usetheme{Madrid}
\usecolortheme{default}
\setbeamertemplate{navigation symbols}{} % Oculta los símbolos de navegación inferiores
\setbeamertemplate{footline}[frame number] % Muestra solo el número de diapositiva
\setbeamercolor{title}{bg=blue!80!black, fg=white}
\setbeamercolor{frametitle}{bg=blue!80!black, fg=white}
\setbeamercolor{itemize item}{fg=blue!80!black}

% Configuración de unidades
\sisetup{
    output-decimal-marker = {.},
    inter-unit-product = \ensuremath{{}\cdot{}}
}

% --- PORTADA ---
\title[IMT-221 \textbar{} Sesión 1]{Circuitos Electrónicos II (IMT-221)}
\subtitle{Bienvenida y Análisis Dinámico de CA}
\author{M.Sc. Ing. Bernardo Quiroga Turdera}
\institute[UCB Tarija]{Universidad Católica Boliviana ``San Pablo'' - Sede Tarija \\ Carrera de Ingeniería Mecatrónica}
\date{03 de agosto de 2026}

\begin{document}

\begin{frame}
    \titlepage
\end{frame}

\begin{frame}{Bienvenidos al Taller de Ingeniería}
    Esta asignatura no es una lista de temas teóricos para memorizar. Es el espacio donde construirán un \textbf{sistema de control analógico de precisión}, real e inmune al ruido industrial.
    
    \vspace{0.5cm}
    \begin{block}{Nuestra Meta}
        Transición sistemática del análisis pasivo de circuitos al control dinámico activo continuo.
    \end{block}
\end{frame}

\begin{frame}{Metodología de Trabajo}
    \begin{itemize}
        \item \textbf{Enfoque ``Python-First'':} Procesamiento de datos experimentales, gráficos de Bode y análisis estadísticos mediante \texttt{numpy}, \texttt{scipy} y \texttt{matplotlib}.
        \vspace{0.3cm}
        \item \textbf{Modelo F.I.V.E. para Laboratorios:}
        \begin{enumerate}
            \item \textbf{F}undamentación (Modelos matemáticos teóricos).
            \item \textbf{I}mplementación (Montaje de hardware real).
            \item \textbf{V}alidación (Contraste entre teoría vs. práctica).
            \item \textbf{E}videncia (Reporte técnico con estándar IEEE).
        \end{enumerate}
    \end{itemize}
\end{frame}

\begin{frame}{Sistema de Evaluación (100\%)}
    \begin{columns}[T]
        \begin{column}{0.5\textwidth}
            \textbf{Evaluación Continua (50\%)}
            \begin{itemize}
                \item \textbf{EC1 (Sem 1-4):} Análisis CA, Impedancias y Bode pasivo. (Incluye Hito 1 PFI).
                \item \textbf{EC2 (Sem 5-17):} Op-Amps reales, Filtros, PID y Potencia. (Incluye Hitos 2 y 3).
                \item \textit{Recuperatorios:} Semanas 5 y 17 para nivelación teórica.
            \end{itemize}
        \end{column}
        \begin{column}{0.5\textwidth}
            \textbf{Evaluación Final Híbrida (50\%)}
            \begin{itemize}
                \item \textbf{Comp. A (40\% Grupal):} Prueba de estrés en vivo al prototipo físico en el laboratorio.
                \item \textbf{Comp. B (60\% Individual):} Examen escrito de rediseño analítico sobre fallas de su propio hardware.
            \end{itemize}
        \end{column}
    \end{columns}
\end{frame}

\begin{frame}{Evaluación Diagnóstica}
    \begin{alertblock}{Instrucciones - Tiempo: 30 minutos}
        Esta prueba no tiene ponderación sumativa, pero es un \textbf{requisito obligatorio} para medir las competencias previas.
    \end{alertblock}
    
    \vspace{0.3cm}
    \textbf{Áreas a evaluar:}
    \begin{itemize}
        \item Álgebra de números complejos y fasores básicos.
        \item Leyes de Kirchhoff en corriente continua.
        \item Dinámica temporal en capacitores.
        \item Parámetros fundamentales de señales senoidales.
    \end{itemize}
    \vspace{0.3cm}
    \centering \textit{(Proceder a la entrega de las hojas de evaluación)}
\end{frame}

\begin{frame}{Tema 1.1: Análisis Dinámico de CA}
    \textbf{Definición Matemática de una Señal Senoidal}
    
    Toda excitación senoidal permanente se modela como:
    \begin{equation*}
        v(t) = V_m \cos(\omega t + \phi)
    \end{equation*}
    
    \begin{itemize}
        \item $V_m$: Amplitud pico de la señal \unit{\volt}.
        \item $\omega = 2\pi f$: Frecuencia angular \unit{\radian\per\second}.
        \item $f = 1/T$: Frecuencia cíclica \unit{\hertz} (\SI{50}{\hertz} nominal en Bolivia).
        \item $\phi$: Ángulo de fase inicial \unit{\radian} o \unit{\degree}.
    \end{itemize}
\end{frame}

\begin{frame}{El Valor Eficaz (RMS)}
    \textbf{Pregunta detonante:} \textit{¿Por qué el multímetro marca \SI{220}{\volt} si la tensión alterna promedia cero en el tiempo?}
    
    \vspace{0.3cm}
    El valor eficaz es el equivalente en CC que entrega la \textbf{misma potencia promedio} a una resistencia pura $R$.
    
    \begin{equation*}
        P_{avg} = \frac{1}{T} \int_{0}^{T} \frac{V_m^2}{R} \cos^2(\omega t + \phi) dt = \frac{V_m^2}{2R} = \frac{V_{RMS}^2}{R}
    \end{equation*}
    
    \begin{equation*}
        \therefore V_{RMS} = \frac{V_m}{\sqrt{2}} \approx 0.7071 \cdot V_m
    \end{equation*}
    
    \begin{exampleblock}{Aplicación de Ingeniería}
        Para la red local de $V_{RMS} = \SI{220}{\volt}$, los componentes deben soportar un pico real de $V_m = 220 \cdot \sqrt{2} \approx \SI{311.12}{\volt}$.
    \end{exampleblock}
\end{frame}

\begin{frame}{Introducción a la Transformación Fasorial}
    Para evitar resolver ecuaciones diferenciales en el dominio del tiempo, transformamos la señal al dominio frecuencial usando la \textbf{Identidad de Euler}:
    
    \begin{equation*}
        e^{\pm j\theta} = \cos(\theta) \pm j\sin(\theta)
    \end{equation*}
    
    El \textbf{Fasor} ($\hat{V}$) captura únicamente la amplitud y la fase, omitiendo la dependencia temporal:
    
    \begin{equation*}
        \hat{V} = V_m e^{j\phi} = V_m \angle \phi
    \end{equation*}
    \begin{equation*}
        \text{O en convención RMS: } \hat{V}_{RMS} = \frac{V_m}{\sqrt{2}} \angle \phi
    \end{equation*}
\end{frame}

\begin{frame}{Próximos Pasos (Laboratorio 1 - Miércoles)}
    \begin{block}{Obligatorio para la próxima sesión (09:00 - 10:45)}
        \begin{enumerate}
            \item \textbf{Software:} Instalar \texttt{LTSpice} y su entorno \texttt{Python (3.10+)}.
            \item \textbf{Lectura:} Capítulo 9 (Análisis sinusoidal) de \textit{Hayt \& Kemmerly}.
        \end{enumerate}
    \end{block}
    
    \vspace{0.3cm}
    \textbf{Anuncio del Laboratorio 1:}
    El miércoles iniciaremos la simulación e instrumentación del \textit{Puente de Wheatstone} alimentado en CA para acondicionar el sensor térmico de nuestro Proyecto Final Integrador.
\end{frame}

\end{document}